\subsection{Introduction \quad / \quad \textthai{บทนำ}}
\begin{paracol}{2}
Quantum mechanics is arguably the most successful theory in history, but it comes with a formidable computational challenge: the "Exponential Wall." In a system with $N$ particles, the size of the \textbf{Hilbert Space}\index{Hilbert Space@ปริภูมิฮิลเบิร์ต (Hilbert Space)} scales as $2^N$, making exact solutions impossible for more than a few dozen particles. Traditional methods like Quantum Monte Carlo (QMC) or Density Functional Theory (DFT) have dominated for decades, but they often struggle with the "sign problem" or high-electron correlation.

Artificial Intelligence offers a new way to represent multi-dimensional wavefunctions. By using Neural Network Quantum States\index{NNQS@NNQS (NNQS)} (NNQS), physicists can now approximate the ground state of many-body systems\index{Quantum Many-Body Problem@ปัญหาหลายอนุภาคเชิงควอนตัม (Quantum Many-Body Problem)} with unprecedented efficiency, effectively "compressing" the quantum complexity.

\switchcolumn

\begin{thai}
กลศาสตร์ควอนตัมเป็นอุปสรรคสำคัญในการคำนวณมาโดยตลอด เนื่องจากปัญหา \textbf{"คำสาปแห่งมิติ (Curse of Dimensionality)"}\index{Curse of Dimensionality@คำสาปแห่งมิติ (Curse of Dimensionality)} เมื่อเราเพิ่มจำนวนอนุภาค พื้นที่สถานะ (State Space) หรือ\textbf{ปริภูมิฮิลเบิร์ต (Hilbert Space)} จะเพิ่มขึ้นแบบเอกซ์โพเนนเชียล ทำให้คอมพิวเตอร์แบบคลาสิกไม่สามารถคำนวณได้โดยตรง ในบทนี้ เราจะสำรวจว่า Neural Network Quantum States (NNQS) ช่วยให้เราสามารถประมาณค่าฟังก์ชันคลื่นสำหรับระบบที่มีอนุภาคจำนวนมากได้อย่างไร
\end{thai}
\end{paracol}
